\section{Object Constraint Programming in Squeak}
\label{sec:ocp_squeak}

%\secmissing{Defining Constraints in Babelsberg/S}
%\secmissing{Constraints over Collections of Objects}
This section describes how constraints are expressed in our Squeak
implementation of Babelsberg, called Babelsberg/S. For our examples,
we use the rules of a Sudoku puzzle.

\begin{lstlisting}[caption=Defining the domain of a Sudoku cell, label=lst:first-Sudoku-constraint]
constraint := [(sudoku at: 1 at: 1) between: 1 and: 9]
                  alwaysSolveWith: solver.
\end{lstlisting}

\autoref{lst:first-Sudoku-constraint} shows the constraint for
defining the domain of one Sudoku cell. In general, a constraint in
Babelsberg/S is specified by sending the message
\texttt{alwaysSolveWith:} to a constraint block. The argument should
be an instance of \texttt{ConstraintSolver}. It is also possible to
solve the constraint with a default constraint solver, which is global
inside the Squeak image, by sending \texttt{alwaysTrue}. As mentioned
in \autoref{sec:introduction}, the constraint is defined by using
object-oriented method definitions rather than a \ac{dsl}. The keyword
\texttt{sudoku} in \autoref{lst:first-Sudoku-constraint} is a variable
of the interactive application and the method \texttt{between:and:} is
a known method in Squeak.

\paragraph{Constraint Construction}
\label{sec:constr-constr}

To construct the constraint, the constraint block is executed in a
different execution mode called \ac{ccm}. Instead of evaluating the
source code normally, in \ac{ccm}, constraint expressions are created
from the constraint code. After that, these constraint expressions are
added to a \texttt{Constraint} object, which is passed to the
constraint solver. We explain the solving process in more detail in
\autoref{sec:maintain_constraints}. If solving succeeded, the method
\texttt{alwaysSolveWith:} returns the newly created constraint
object. This object can then be used to reflect on the constraint
(e.g., to inspect which variables participate in it), as well as to
dynamically disable and re-enable the constraint later. If solving
fails, an exception is raised which must be handled by the
programmer. In that case, the constraint is not added and the system
remains unperturbed.

Squeak/Smalltalk includes an in-image Smalltalk interpreter. Squeak
stack frames can be reified into subinstances of the
\lstinline!ContextPart! class. These provide methods to interpret each
bytecode. This facility is used by the Squeak debugger.

Babelsberg/S uses a modified version of the image-side interpreter to
evaluate the constraint block. The \lstinline!alwaysTrue! method
creates a new Process (a Smalltalk green-thread) from the constraint
block. This process is interpreted stepwise using the interface of the
ContextPart objects. Where interpretation in \ac{ccm} deviates from
normal Smalltalk semantics, we use
ContextS\cite{contextShirschfeld2008introduction} to layer methods
whose behavior needs to change inside a \ac{ccm} layer.

We modified the interpreter to create \texttt{ConstraintVariable}
objects for instance variables that are accessed through accessor
methods. We wrap the accessors to go through these variables to read
and write values. All methods, including operator methods such as
\textbf{+} and \textbf{-}, are then called on constraint variables and
construct constraint expressions instead of evaluating directly.

Back to the constraints of a Sudoku puzzle, we must define the domain
of all cells inside the \texttt{sudoku}. The usual way in Babelsberg
is shown in \autoref{lst:loop-Sudoku-constraint}. It uses a loop over
all Sudoku cells.

\begin{lstlisting}[caption=Defining the domain of all Sudoku cells with a loop, label=lst:loop-Sudoku-constraint]
(1 to: sudoku size) do: [:index | 
		[(sudoku at: index) between: 1 and: 9] 
				alwaysSolveWith: solver].
\end{lstlisting}

Babelsberg/Squeak contributes to more human-readable constraints by supporting the Collection \ac{api} in constraints. As a result, the domain constraint of a Sudoku puzzle can be expressed through sending \texttt{allSatisfy:} to \texttt{sudoku}, as shown in \autoref{lst:collectionAPI-Sudoku-constraint}.

\begin{lstlisting}[caption=Defining the domain of all Sudoku cells with the Collection \ac{api}, label=lst:collectionAPI-Sudoku-constraint]
[sudoku allSatisfy: [:cell | cell between: 1 and: 9]] 
		alwaysSolveWith: solver.
\end{lstlisting}

The extension of Babelsberg/S for the Collection \ac{api} slightly changes the implementation. For maintaining the constraints automatically, as explained in \autoref{sec:maintain_constraints}, read and write access on constraint variables, e.g. on Sudoku cells, has to be detected. In Squeak, an access on an arbitrary collection calls \texttt{at:} or \texttt{at:put:} on an internal array, that maintains the elements of the collection. Therefore, sending the messages \texttt{at:} and \texttt{at:put:} to an array have to be modified. To reach this, in \ac{ccm}, the internal array is replaced by a \texttt{ConstraintArray}, that inherits from \texttt{Array} and overrides the accessor methods \texttt{at:} and \texttt{at:put:}. 

%%% Local Variables:
%%% mode: latex
%%% TeX-master: "paper"
%%% End:
